\section{The same Fubini--Study function in arbitrarily small perturbations of its class}
\label{sec:native-numerical-psh-power-perturbations}

We keep the notation of Section~\ref{sec:native-numerical-psh-generated-twists}.

\begin{lemma}[Powering a varying-line datum]
\label{lem:native-numerical-fs-power-datum}
\lean{EnvelopeProperty.NativeNumericalFubiniStudy.FSData.powerDatum,EnvelopeProperty.NativeNumericalFubiniStudy.FSData.powerDatum_eval,EnvelopeProperty.NativeNumericalFubiniStudy.FSData.powerDatum_numericalClass,EnvelopeProperty.NativeNumericalFubiniStudy.FSData.IsRational.powerDatum,EnvelopeProperty.NativeNumericalFubiniStudy.FSData.Nonzero.powerDatum}
\leanok
\uses{def:native-numerical-whole-section-fs-data,def:scheme-global-section-tensor-power,thm:scheme-global-section-tensor-power-norm,lem:native-numerical-line-tensor-power,thm:native-numerical-fs-nonvanishing}
\source{bj22-trivval:equ:FSpower}
Let $D=(M,m,(s_j,w_j)_j)$ be a varying-line Fubini--Study datum and $r>0$ an
integer.  The datum $D^{(r)}=(M^{\otimes r},mr,(s_j^{\otimes r},rw_j)_j)$ has
the same function and the same class as $D$.  It has rational weights if $D$
does, and, if $X$ is integral, it is nonzero if $D$ is.
\end{lemma}
\begin{proof}
The norm of $s_j^{\otimes r}$ at a point is the $r$-th power of the norm of
$s_j$, so $\log\lVert s_j^{\otimes r}\rVert+rw_j=r(\log\lVert s_j\rVert+w_j)$,
also when $s_j$ vanishes at the point.  Multiplication by $r>0$ commutes
with finite maxima in $[-\infty,+\infty)$, and $(mr)^{-1}r=m^{-1}$; hence
the functions agree.  The class of $M^{\otimes r}$ is $r[M]$, so
$[D^{(r)}]=(mr)^{-1}r[M]=[D]$.  The weights $rw_j$ are rational when the
$w_j$ are.  Nonzeroness follows from the equality of functions on an
integral scheme (Theorem~\ref{thm:native-numerical-fs-nonvanishing}).
\end{proof}

\begin{theorem}[The same function in arbitrarily small generated perturbations]
\label{thm:native-numerical-fs-small-generated-perturbations}
\lean{EnvelopeProperty.NativeNumericalFubiniStudy.FSData.exists_small_generatedPerturbations}
\leanok
\uses{lem:native-numerical-fs-power-datum,thm:native-numerical-generated-twist-exact,thm:native-numerical-curve-weak-vector-topology}
Let $X$ be integral and projective, let $D$ be a nonzero varying-line
Fubini--Study datum with denominator $m$, and let $T$ be an invertible sheaf
with a finite family of global generators.  There are nonzero data $E_j$,
$j\in\mathbf N$, such that at every point $y$
\[
 \varphi_{E_j}(y)=\varphi_D(y),\qquad
 [E_j]=[D]+\frac{1}{m(j+1)}[T],\qquad
 [E_j]\longrightarrow[D].
\]
\end{theorem}
\begin{proof}
Let $E_j$ be the generated twist by $T$
(Theorem~\ref{thm:native-numerical-generated-twist-exact}) of the power
$D^{(j+1)}$ (Lemma~\ref{lem:native-numerical-fs-power-datum}).  Both
operations preserve the function and nonzeroness; the power keeps the class
and has denominator $m(j+1)$, and the twist adds $(m(j+1))^{-1}[T]$.  Since
$(m(j+1))^{-1}\to0$ and addition and scalar multiplication are continuous
for the curve-degree topology, $[E_j]\to[D]$.
\end{proof}

\begin{corollary}[Perturbations in an ample direction]
\label{cor:native-numerical-fs-small-ample-direction}
\lean{EnvelopeProperty.NativeNumericalFubiniStudy.FSData.exists_small_ampleDirectionPerturbations}
\leanok
\uses{thm:native-numerical-fs-small-generated-perturbations,lem:scheme-ample-line-eventual-generation,lem:native-numerical-line-tensor-power}
\source{bj22-trivval:lem:twopsh,bj22-trivval:thm:pshample}
Let $X$, $D$ and $m$ be as in
Theorem~\ref{thm:native-numerical-fs-small-generated-perturbations} and let
$A$ be an ample invertible sheaf.  There are an integer $q>0$ and nonzero
data $E_j$, $j\in\mathbf N$, such that at every point $y$
\[
 \varphi_{E_j}(y)=\varphi_D(y),\qquad
 [E_j]=[D]+\frac{q}{m(j+1)}[A],\qquad
 [E_j]\longrightarrow[D].
\]
\end{corollary}
\begin{proof}
Some positive power $A^{\otimes q}$ has a finite family of global generators
(Lemma~\ref{lem:scheme-ample-line-eventual-generation}).  Apply
Theorem~\ref{thm:native-numerical-fs-small-generated-perturbations} to
$T=A^{\otimes q}$ and use $[A^{\otimes q}]=q[A]$.
\end{proof}

\begin{remark}
The line of $E_j$ is $M^{\otimes(j+1)}\otimes A^{\otimes q}$.  No ampleness of
this line is asserted.  In BJ22 these perturbations are combined with the
openness of the ample cone of the finite dimensional space $N^1(X)$, with
Kleiman's criterion (ampleness depends only on the numerical class), and
with Proposition~2.19(iii) for an ample difference line, in order to place
$\varphi_D$ in $\mathcal H(L)$ for a $\mathbb Q$-line bundle $L$ of a
prescribed class.
\end{remark}
